\chapter{Introduction}
\label{ch:introduction}
\resetcounterformats

% The introduction motivates all three papers, sets out the contribution to
% the literature, and gives the reader a map of the thesis.
% DRAFT 13 Sep 2026: Motivation and Contribution written from the current
% chapters; Overview rewritten to match the chapters' present findings.
% For Damian's review and markup.

\section{Motivation}
\label{sec:intro_motivation}

Schools with similar pupils and similar budgets produce very different
results. Part of the difference comes from what the schools do, not from
which pupils they take. This thesis measures two of the things a school
does: the culture of the school, and whether it offers a high-return
subject. Culture is hard to measure, and the subject offer is rarely
measured. The thesis measures both for English secondary schools, and asks
what each one means for pupils.

The first two papers are about school culture. Nearly every school describes itself
as warm, supportive and well ordered. The measures available to researchers
and to policy do not settle whether that is true. An inspection grade
reflects exam results as well as culture, so a good grade does not tell us
whether a school has a good culture. A high exclusion rate could mean a school that enforces its rules
firmly or one that has lost control, and the two cannot be told apart from
the rate alone. Surveys record what people say rather than what a visitor would
see.

The first two papers measure culture by observation instead. A research team that I led
visited \NTierOne{} schools, spent a structured day observing each one, and scored
two dimensions of its culture: warmth, the relational quality of the school,
and strictness, its behavioural structure. The two dimensions come from the
parenting literature, where their combination --- demanding and responsive
at once --- is the style that produces the best child outcomes. Separating
them matters because the public debate about schools tends to treat them as
opposites: the case for strict schools is often put as a case against soft
ones. Whether warmth and strictness are substitutes, complements or simply
independent is an empirical question, and answering it requires measuring
each on its own.

Direct observation produces credible scores, but it cannot be done for
every school. Every school, however, produces text: an
inspection report written about it, and a behaviour policy and a website
written by it. The first paper asks how much of what an observer sees in a school
can be predicted from the text that exists for every school, and the second asks
what the observed scores predict. The two questions depend on each other.
The visits give a direct measure of each school's culture. The first paper
checks the text-based scores against that direct measure. Only the scores
from inspection reports agree with the visit scores. The second paper
therefore uses the visit scores for the visited schools, and the
inspection-report scores when it extends the analysis to every school in
England. Scores from behaviour policies and websites are not used.

The third paper is about the subject a school offers. It is simpler to measure than
culture. A school either offers the subject or it does not. Whether a pupil can take A-level
economics --- the usual route into one of the highest-earning degrees
--- depends on whether their school employs a teacher qualified to teach
it, and during the period we study fewer than half of English state
schools did. If a school does not offer the subject, its pupils cannot take
it. The
third paper measures what this constraint costs in lifetime earnings,
establishes who it falls on, and asks whether the tax revenue from removing
it would cover the cost of doing so.

The three papers share a method as well as a subject. Each measures
something about schools for all, or nearly all, English secondary schools:
the culture a school practises, how much of that culture appears in the
documents written by and about the school, and whether a school offers a
subject. The method combines direct measurement with
administrative data. In the first two papers it also uses language models,
whose scores are checked against the visits. The thesis shows that these
features of schools can be measured for every school, and reports what the
measurements show.

\section{Contribution to the Literature}
\label{sec:intro_contribution}

The first paper contributes to two literatures. The school-climate
literature has long used the parenting typology of demandingness and
responsiveness to classify schools, and finds that pupils do best in
schools that are both demanding and responsive; but it measures climate
through surveys of pupils, parents or staff, which record perceptions
rather than practice. The paper's first contribution is a gold standard
built instead from direct observation. It draws an explicit distinction
between what a school enacts and what its headteacher espouses, which the
survey literature cannot draw. The distinction turns out to be large:
headteachers' accounts of their schools differ from what observers record,
and the difference is larger for warmth than for strictness.

The second contribution is to the
text-as-data literature: the paper uses the documents that exist for every
school as a source of measurement, scores them with a prediction model and
with a language model working to a mark scheme, and validates both against
the gold standard. The paper finds which documents contain information
about a school's warmth and strictness, and which do not. Inspection
reports contain information about both. Warmth is predicted by the words a report ordinarily uses.
Strictness is found only in specific statements, such as that conduct
standards are consistently enforced, which the model must be told to look
for. Behaviour policies and websites contain no usable information about
either. Their scores agree with what the headteacher says about
strictness, and with nothing about warmth. Of the three documents that
exist for every school, only the one written by an inspector says anything
about its culture.

The second paper contributes to the literature on school effectiveness. The
most credible causal evidence on school culture comes from studies of
high-expectation charter schools in the United States. Those schools
combine several practices at once, such as strict discipline, high
expectations, extended time and intensive tutoring, so the studies
estimate the effect of the whole package. This paper separates warmth from
strictness and estimates each one's association with progress. In the
visited schools the two scores are correlated at only about 0.36, so a
school's warmth can be distinguished from its strictness. Both are
associated with progress: a standard deviation of warmth with
\PrimarySpecWSD{} of a standard deviation of Progress~8, and a standard
deviation of strictness with \PrimarySpecSSD{}. Schools in the top half on
both dimensions make more progress than the rest, and their advantage is
the sum of the two separate associations. Having both does not add
anything beyond that sum. Controlling for the
quality of the teaching observed on the visit does not reduce either
association.

These estimates are associations, not causal effects. Schools that differ in
culture also differ in other ways that are not observed: a headteacher who
builds a warm culture may also be good at other things that raise
progress, and parents who value a warm school may choose one for children
who would progress well anyway. The estimates therefore show that pupils
make more progress at warmer and stricter schools, not that changing a
school's culture would change its pupils' progress. The paper's
contribution is to measure culture by observation and to describe how it
relates to progress. Warmth receives less attention in policy
debate than strictness, but it predicts progress just as strongly, and this
is only visible when warmth is measured by observation rather than
self-report.

The third paper, joint with Arun Advani and Elaine Drayton, contributes to
the literature on field-of-study returns. That literature measures the
value of a subject by comparing pupils just above and below an entry
cut-off. Our contribution is to measure it where the question is not who
wins a place but whether the subject exists at the pupil's school at all.
We use linked administrative data covering every pupil in England. Schools
gain and lose the subject when a qualified teacher arrives or leaves, and
we compare pupils at the same school in the years before and after such a
change. We show that pupils at schools without the subject do not take it,
that those who miss it go on to lower-paying degrees and jobs, and that
the schools without the subject are concentrated outside London and the
South East.

The paper then costs the policy. One teacher is enough to let a
school offer the subject. Pupils who take the subject go on to earn more,
and so pay more tax over their careers. We compare that extra tax with the
cost of employing the teacher.
The return exceeds the cost several times over.

Together the papers make one argument in two settings. What a school does
--- the culture it enacts day to day, the subjects it puts on its timetable
--- is measurable, varies widely between otherwise similar schools, and
predicts what its pupils go on to achieve.

\section{Overview of the Thesis}
\label{sec:intro_overview}

The remainder of this thesis is organised as follows.
\Cref{ch:paper1} develops the measurement framework. It constructs
gold-standard warmth and strictness scores for \NTierOne{} schools from
school visits and headteacher interviews. It then asks how far those scores
can be predicted from text that exists for every school: Ofsted inspection
reports, behaviour policies and school websites. Two methods are used. A
prediction model learns which words and phrases go with warmer or stricter
schools, and a large language model scores each document against a mark
scheme. Inspection reports predict both warmth and strictness, though only
modestly, and in different ways. Scores from behaviour policies and websites match what headteachers
say about their schools. They do not match what observers record on the
visit.

\Cref{ch:paper2} uses the observed scores as explanatory variables. It
estimates the association between school culture and the progress pupils
make over secondary school, measured by Progress~8, the Department for
Education's official progress measure. Both warmth and strictness predict
progress. One rival explanation is that warm or strict schools simply have
better teaching, and that the association is really with teaching quality.
Adding the quality of the teaching observed on the visit as a control
leaves both associations unchanged, so that is not what drives them. Warmth predicts progress only when
it is measured by observation: replacing the observers' scores with the
headteachers' own ratings removes the warmth association, whereas
strictness is associated with progress under either measure. The text-based
scores from inspection reports exist for nearly every English state
secondary school. Estimating the same models with those scores, across all
English state secondary schools, gives the same result: both warmth and
strictness predict progress.

\Cref{ch:paper3}, joint work with Arun Advani and Elaine Drayton, turns
from culture to curriculum. It shows that whether an English school offers
A-level economics depends on whether it employs a suitably qualified
teacher, compares pupils at the same school before and after the subject
arrives or disappears to estimate the effect of taking the subject on
degree choice and earnings, and computes the fiscal
return to expanding the supply of economics teachers, which exceeds the
cost of a teacher several times over.
